\documentclass[main]{subfiles}


\title[Advanced Hyperparameter Optimization]{Advanced Hyperparameter Optimization}
\subtitle{Multi-Fidelity Optimization}
\author[Katharina Eggensperger]{Katharina Eggensperger}
\date{}

\titlebackground*{images/AutoML_Speeding_up_HPO.jpg}

%\institute{}
%\date{}

\begin{document}
	
	\maketitle
	

%----------------------------------------------------------------------
%----------------------------------------------------------------------
\begin{frame}[c]{Should we always completely evaluate a pipeline?}

$\leadsto$ Full evaluations yield the most accurate signal, but may waste resources in practice.

\vspace{0.5cm}

In practice, we often have access to cheap approximations or intermediate results of our actual target function:
\begin{itemize}
    \item subsampling training data
    \item partial training
    \item downscaled images
    \item smaller models
\end{itemize}

\vspace{1cm}
\textbf{$\leadsto$ How can we use this information to increase the efficiency of Bayesian optimization?} \\

\end{frame}


%-----------------------------------------------------------------------
%----------------------------------------------------------------------
\begin{frame}[c]{A motivating example}

\begin{center}
\includegraphics[width=0.8\textwidth]{05_Advanced_HPO/images/multi-fidelity_example.png}
\end{center}
\liturl{Franceschi et al. "Hyperparameter Optimization in Machine Learning"}{https://arxiv.org/abs/2410.22854}

\end{frame}


%----------------------------------------------------------------------
%----------------------------------------------------------------------
\begin{frame}[c]{What are intermediate results, early metrics or proxy performance?}

Evaluating ML systems is often iterative and we can get intermediate results such as:
\begin{itemize}
    \item subsample dataset
    \item partial training
    \item coarse simulation / downscaled images
    \item smaller model
\end{itemize}

$\leadsto$ These \textit{early} metrics are not necessarily reliable, but might be useful to identify promising regions and discard bad-performing candidates early.
$\leadsto$ Optimization methods that make use of these are called multi-fidelity methods.

\end{frame}
%-----------------------------------------------------------------------
%----------------------------------------------------------------------
\begin{frame}[c]{Formal definition}

We refer to an inexpensive approximation of the true objective $f(\lambda)$ as $f_b(\lambda)$, where $b \in [b_{min}, b_{max}]$ with $c(\lambda,b)$ being its associated cost. Furthermore, we assume that \item $f_{b_{max}}(\lambda) = f(\lambda)$.

\\~\\
Our optimization objective is then to find:

\begin{equation}
\lambda^* \in \argmin_{\lambda \in \Lambda} f_{b_{max}}(\lambda)
\end{equation}

\end{frame}
%-----------------------------------------------------------------------
%----------------------------------------------------------------------
\begin{frame}[c]{Assumptions}

These methods exploit three structural observations
\begin{enumerate}
    \item the \textbf{error shrinks smoothly} as the fidelity budget increases,
    \item \textbf{ranks are approximately consistent} across fidelities,
    \item and \textbf{higher fidelities provide a more accurate approximation} of the true objective value.
\end{enumerate}

$\leadsto$ in practice, these assumptions are not necessarily satisfied. HPO must be robust against that.

\end{frame}
%-----------------------------------------------------------------------
%----------------------------------------------------------------------
\begin{frame}[c]{Optimization target}

<add multi-fidelity optimization target here>

Note that this is also studied in Bayesian optimization with environmental variables (double check), e.g. engineering and material design. 

\end{frame}
%-----------------------------------------------------------------------
%----------------------------------------------------------------------
\begin{frame}[c]{Successive Halving (add citation here)}

<add multi-fidelity optimization target here>

textual explanation of the algorithm
\end{frame}
%-----------------------------------------------------------------------
%----------------------------------------------------------------------
\begin{frame}[c]{Successive Halving Details}

add sth about guarantees and what they mean; mention failure cases

\end{frame}
%-----------------------------------------------------------------------
%----------------------------------------------------------------------
\begin{frame}[c]{Hyperband (add citation here)}

textual description

\end{frame}
%-----------------------------------------------------------------------
%----------------------------------------------------------------------
\begin{frame}[c]{Model-based multi-fidelity optimization}

Can we combine this with Bayesian optimization? - BOHB


\end{frame}
%-----------------------------------------------------------------------
%----------------------------------------------------------------------
\begin{frame}[c]{Can we do better?}

Talk about learning curve prediction, freeze-thaw, etc.

(note to myself: Should we mention bandits in this presentation? Probably not.)

\end{frame}
%-----------------------------------------------------------------------
%----------------------------------------------------------------------
\begin{frame}[c]{What if the assumptions not apply?}


\end{frame}
%-----------------------------------------------------------------------
%----------------------------------------------------------------------
\begin{frame}[c]{Summary}


\end{frame}
%-----------------------------------------------------------------------
%----------------------------------------------------------------------

\end{document}